\section*{术语索引}
\phantomsection
\addcontentsline{toc}{section}{术语索引}

参考编号依次表示章、节、小节（或习题）。（R 指“结果概要”。）

\begin{Shaded}
\begin{Highlighting}[]
\NormalTok{对有序集添加最大元素：III.1.7}
\NormalTok{添加（将一个集合添加到另一个集合）：R.5.4}
\NormalTok{在集合上一致（函数）：II.3.5}
\NormalTok{（两个函数在集合上）一致：R.2.14}
\NormalTok{阿列夫：III.6. 习题 10}
\NormalTok{前件组合式：I. 附录 4}
\NormalTok{反定向（有序集）：III.1. 习题 23}
\NormalTok{将一个理论的结果应用于另一个理论：I.2.4}
\NormalTok{自变量：R.1.2}
\NormalTok{组合式：I.1.1}
\NormalTok{前件：I. 附录 4}
\NormalTok{平衡的：I. 附录 4}
\NormalTok{第一种（第二种）的：I.1.3}
\NormalTok{完全平衡的：I. 附录 4}
\NormalTok{相伴等价关系：II.6.2}
\NormalTok{积结构的结合性判据：IV.2.4}
\NormalTok{集合族的交的结合性：R.4.8}
\NormalTok{集合族的积的结合性：R.4.11}
\NormalTok{两个集合的并与交的结合性：R.1.14}
\NormalTok{集合族的并的结合性：R.4.3}
\NormalTok{自同构：IV.1.5，R.8.6}
\NormalTok{辅助基集：IV.1.3，1.4}
\NormalTok{常数，方法：I.3.3}
\NormalTok{假设，方法：I.3.3}
\NormalTok{公理，显式：I.2.1}
\NormalTok{隐式：I.2.1}
\NormalTok{选择公理：R.4.10}
\NormalTok{外延公理：II.1.3}
\NormalTok{无穷公理：III.6.1}
\end{Highlighting}
\end{Shaded}

\begin{Shaded}
\begin{Highlighting}[]
\NormalTok{有序对公理：II.2.1}
\NormalTok{子集集合公理：II.5.1}
\NormalTok{二元集合公理：II.1.5}
\NormalTok{公理，等价的：R.8.4}
\NormalTok{同种结构公理：R.8.2}
\NormalTok{结构种公理：IV.1.4}

\NormalTok{平衡组合：I. 附录 4}
\NormalTok{平衡词：I. 附录 2}
\NormalTok{展开的基，或记数系统的基：III.5.7}
\NormalTok{集合标度的基：R.8.1}
\NormalTok{集合（层级构造中的）基：IV.1.1}
\NormalTok{辅助基集合：IV.1.3, 1.4}
\NormalTok{主基集合：IV.1.3, 1.4}
\NormalTok{双射：II.3.7, R.2.9}
\NormalTok{双射映射：II.3.7, R.2.9}
\NormalTok{二项式系数：III.5.8}
\NormalTok{布尔格：III.1. 习题 17}
\NormalTok{界，最大下界和最小上界（集合或映射的）：III.1.9, R.6.7}
\NormalTok{最小上界（基数集合的）：III.3.2}
\NormalTok{下界：III.1.8, R.6.7}
\NormalTok{严格上界：III.2.4}
\NormalTok{上界：III.1.8, R.6.7}
\NormalTok{有界的，上有界的，下有界的（集合或映射）：III.1.8, R.6.7}
\NormalTok{分支的（有序集合）：III.1. 习题 24}
\end{Highlighting}
\end{Shaded}

函数的典范分解：II.6.5, R.5.3 典范扩张，对应关系到子集集合的：II.5.1 函数族到乘积集合的：II.5.7 映射的：IV.1.2 两个函数到乘积集合的：II.3.9 带符号的：IV.2. 习题 1 典范注入：II.3.7 典范映射，IV, 附录 4 从E的子集到E的：II.3.7, R.2.3 从A的\(^{B \times C}\) 到 (A\(^{B}\))\(^{C}\) : II.5.2, R.4.14 将 G 映到 G (G 为图) : II.3.7 将 F\(^{E}\) 映到 F(E, F) : II.5.2 将 F(B × C, A) 映到 F(B, F(C, A)) : II.5.2 将 ∏\emph{\{i∈\{α\}\} X\_i 映到 X\_α : II.5.3 将 ∏}\{i∈\{α, β\}\} X\_i 映到 X\_α × X\_β : II.5.3

典范映射 \(\prod_{i\in I}X_t\) 到 \(\prod_{\lambda \in L}\left(\prod_{i\in J_\lambda}X_t\right)\) 上的: II.5.5, R.4.11\\
\(\prod_{i\in I}X_t\) 到 \(\left(\prod_{i\in J_\alpha}X_t\right) \times \left(\prod_{i\in J_\beta}X_t\right)\) 上的 ((J\(_\alpha\), J\(_\beta\)) 为 I 的一个划分): II.5.5\\
\(\prod_{i\in \{\alpha,\beta,\gamma\}}X_t\) 到 \(X_\alpha \times X_\beta \times X_\gamma\) 上的: II.5.5\\
\(\prod_{i\in I}X_t^E\) 到 \(\left(\prod_{i\in I}X_t\right)^E\) 上的: II.5.7, R.4.13\\
将 E 映到 E/R: II.6.2, R.5.2\\
将 A/R\(_A\) 映到 f\textless A\textgreater: II.6.6, R.5.5\\
将 (E/S)/(R/S) 映到 E/R: II.6.7, R.5.9\\
将 (E × E')/(R × R') 映到 (E/R) × (E'/R'): II.6.8, R.5.10\\
到由指标集限制得到的直接极限的直接极限: III.7.6\\
将 E\(_\beta\) 映到 lim E\(_\alpha\): III.7.5, R.6.13\\
将逆极限映到由指标集限制得到的逆极限: III.7.1\\
将 lim E\(_\alpha\) 映到 E\(_\beta\): III.7.1, R.6.14\\
将 E × F 映到 F × E: R.3.4\\
将 E × F × G 映到 (E × F) × G 等: R.3.12\\
将 FA 映到 \(\prod_{i\in I}FA_i\) (其中 A = \(\bigcup_{i\in I}A_i\), 且当 t ≠ x 时 A\(_t\) ∩ Ax = φ): R.4.15\\
对称性: R.3.4\\
康托尔定理: III.3.6\\
基数，主导的: III.6.Ex.21\\
有限的: III.4.1\\
不可达的: III.6.Ex.22\\
积: III.3.3\\
正则的: III.6.Ex.17\\
奇异的: III.6.Ex.17\\
强不可达的: III.6.Ex.22\\
集合的基数: III.3.1\\
基数和: III.3.3\\
有序集中元素的链（关于一个映射）: III.2.Ex.6\\
集合的子集的特征函数: III.5.5\\
特征刻画，典型的: IV.1.4\\
选择公理: R.4.9\\
类，等价: II.6.2, R.5.2\\
与 x 等价的对象类（关于一个等价关系）: II.6.9\\
闭区间: III.1.13, R.6.4\\
闭包（有序集到自身的映射）: III.1.Ex.13

较粗的等价关系 : II.6.7 预序 : III.1.4 结构 : IV.2.2 系数，二项式 : III.5.8 共尾子集 : III.1.7, R.6.5 重合（两个函数在集合上） : II.3.5, R.2.14 共首子集 : III.1.7, R.6.5 集合化关系 : II.1.4 交换性（并和交的） : R.1.14 可比较的元素 : III.1.12 结构 : IV.2.2 相容性，函数与两个等价关系的 : II.6.5, R.5.8 关系与等价关系的 : II.6.5, R.5.7 集合的补集 : II.1.7, R.1.7 完全格 : III.1.Ex.11 完全解 : I.5.2 完全分歧的（有序集） : III.2.Ex.8 有序集的完备化 : III.1.Ex.15 复合，映射的 : R.2.11 集合的 : R.3.10 两个对应的 : II.3.3 两个图的 : II.3.3 两个关系的合取 : I.3.4 集合关于关系的连通分量 : II.6.Ex.10 常项，辅助的 : I.3.3 函数或映射 : II.3.4, R.2.3 理论的 : I.2.1 集合关于关系的构成部分 : II.6.Ex.11 构造，层级 : IV.1.1 形成的 : I.1.3 包含于集合中 : II.1.2, R.1.12 连续统假设 : III.6.4 连续统，势 : III.6.4 矛盾的公理 : R.8.6 理论 : I.2.2 逆变有符号层级类型 : IV.2.Ex.1 坐标（第一，第二）有序对的 : II.2.1 坐标函数（第一，第二） : II.3.6 指标 i 的 : II.5.3 坐标函数，集合族乘积上的 : R.4.11 若干集合乘积上的 : R.3.12 两个集合乘积上的 : R.3.1 推论 : I.2.2 对应，两个集合之间的 : II.3.1

在对象 x 处定义的对应：II.3.1 由关系定义的：II.3.1 逆的：II.3.2 一一的：II.3.7，R.2.9 对应，复合：II.3.3 可数集：III.6.4，R.7.7 可数并、交、积：R.7.9 协变带符号的层式类型：IV.2.Ex.1 覆盖：II.4.6，R.4.4 加细的：II.4.6 集合的覆盖：II.4.6 准则，演绎的：I.2.2 构成的：I.1.4 演绎的准则：I.3.3 代入的准则：I.1.2 临界序数：III.6.Ex.13

合式集合：III.2.Ex.20 十进制系统：III.5.7 分解，典范的：II.6.5，R.5.3 递减子集族：III.1.5，R.6.12 递减映射：III.1.5，R.6.12 导出的（结构）：IV.1.6 演绎，准则：I.3.3 程序：IV.1.6 演绎准则：I.2.2 定义：I.1.1 覆盖的不相交度：III.6.Ex.25 论证文本：I.2.2 递减归纳：III.4.3 对角线，A × A 的：II.3.3，R.3.4 E¹ 的：II.5.3 对角映射，A 到 A × A 的：II.3.7，R.3.4 E 到 E¹ 的：II.5.3 图：R.2.2 两个整数之差：III.5.2 不同于：I.5.1 数字：III.5.7 结构的直接像：IV.2.6 直接极限，映射直接系的：III.7.6，R.6.13 集合直接系的：III.7.5，R.6.13 直接系，映射的：III.7.6，R.6.13 集合的：III.7.5，R.6.13 子集的：III.7.6

有向的（左、右）：III.1.10，R.6.8 关于关系 ≤ 的：III.1.10 不相交区间：I.App.1 集合：II.4.7，R.1.13 析取，情形的方法：I.3.3 两个关系的析取：I.1.3 分配格：III.1.Ex.16 分配律，集合族的并与交的：R.4.3，4.8 两个集合的并与交的：R.1.14 发散映射（序数到自身的）：III.6.Ex.23 可被整数整除：III.5.6 除法，欧几里得：III.5.6 整数的除法：III.5.6 定义域，对应的：II.3.1 图的：II.3.1 优势基数：III.6.Ex.21 双重族：II.3.4 双重序列：III.6.1，R.7.8 对偶公式：R.1.15，R.4.7 对偶规则：R.1.15，R.4.7 二进制系统：III.5.7

层式，层式构造，层式构造方案：IV.1.1 带符号的：IV.2.Ex.1 类型：IV.1.Ex.1 元素：II.1.1 固定的（在映射或映射集下）：R.2.3 一般的：R.1.2 最大的：III.1.3，R.6.5 不变的（在映射或映射集下）：R.2.3 最小的：III.1.3，R.6.5 极大的：III.1.6，R.6.6 极小的：III.1.6，R.6.6 元素，可比较的：III.1.12 空函数：II.3.4 空集：II.1.7，R.1.8 空字：I.App.1 具有结构的：IV.1.4 区间的端点：III.1.13，R.6.4 相等：I.5.1 等词理论：I.5.1 相等，关系：I.5.1，R.1.6 方程：I.5.2 等势集合：III.3.1，R.7.1

\begin{Shaded}
\begin{Highlighting}[]
\NormalTok{等价，在集合上：II.6.1}
\NormalTok{等价类：II.6.2，R.5.2}
\NormalTok{等价关系：II.6.1，R.5.2}
\NormalTok{与函数相关联的：II.6.2}
\NormalTok{较粗的：II.6.7}
\NormalTok{较细的：II.6.7}
\NormalTok{诱导的：II.6.6}
\NormalTok{在集合上：II.6.1}
\NormalTok{积：II.6.8}
\NormalTok{商：II.6.7}
\NormalTok{等价公理：R.8.4}
\NormalTok{等价元素：II.6.1}
\NormalTok{等价势：R.7.2}
\NormalTok{等价关系：I.3.5，R.1.3}
\NormalTok{等价结构种：IV.1.7}
\NormalTok{等价理论：I.2.4}
\NormalTok{欧几里得除法：III.5.6}
\NormalTok{偶整数：III.5.6}
\NormalTok{存在量词：I.4.1}
\NormalTok{整数按底 b 的展开：III.5.7}
\NormalTok{显式公理：I.2.1}
\NormalTok{幂运算：R.4.9}
\NormalTok{扩张，典范的：IV.1.2}
\NormalTok{逆（映射到子集集合的）：R.2.6}
\NormalTok{对应或映射到子集集合的扩张：II.5.1，R.2.4}
\NormalTok{函数族到积集合的扩张：II.5.7}
\NormalTok{映射的扩张：II.3.5，R.2.13}
\NormalTok{序、序关系、预序、预序关系的扩张：III.1.4，R.6.1}
\NormalTok{若干映射到积集合的扩张：R.3.14}
\NormalTok{两个函数到积集合的扩张：II.3.9}
\NormalTok{外延，公理：II.1.3}

\NormalTok{因子，整数的：III.5.6}
\NormalTok{积的：II.2.2 和 II.5.3}
\NormalTok{因子，积集合的：R.3.1，4.9}
\NormalTok{阶乘 n：III.5.8}
\NormalTok{分解，典范的：R.5.3}
\NormalTok{映射的分解：R.2.11}
\NormalTok{假关系：I.2.2}
\NormalTok{族：II.3.4}
\NormalTok{双重：II.3.4}
\NormalTok{有限的: III.4.1}
\NormalTok{集合元素的族：II.3.4，R.2.14}
\NormalTok{两两不相交集合的族：II.4.7}
\end{Highlighting}
\end{Shaded}

\begin{Shaded}
\begin{Highlighting}[]
\NormalTok{子集族，递减（递增）：III.1.5}
\NormalTok{集合的子集族：II.3.4}
\NormalTok{终特征（全序集合的）：III.6.Ex.16}
\NormalTok{字的终段：I.App.1}
\NormalTok{终结构：IV.2.5}
\NormalTok{加细覆盖：II.4.6}
\NormalTok{加细预序：III.1.4}
\NormalTok{加细关系：II.6.7}
\NormalTok{加细结构：IV.2.2}
\NormalTok{有限基数：III.4.1}
\NormalTok{有限特征（性质，子集集合的）：III.4.5，R.6.11}
\NormalTok{有限族：III.4.1}
\NormalTok{有限序列，集合元素的有限序列：III.5.4，R.7.8}
\NormalTok{有限集合：IV.4.1}
\NormalTok{有序对的第一坐标：II.2.1}
\NormalTok{图的第一投影：II.3.1}
\NormalTok{第一类：I.1.3}
\NormalTok{序列的第一项：III.5.4}
\NormalTok{固定元素：II.3.4，R.2.3}
\NormalTok{构造性构造：I.1.3}
\NormalTok{准则：I.1.4}
\NormalTok{有序集的自由子集：III.1.Ex.5}
\NormalTok{函数：见映射}
\NormalTok{特征函数：III.5.5}
\NormalTok{坐标：II.3.6，II.5.3}
\NormalTok{定义在A上取值于B中：II.3.4}
\NormalTok{不依赖于x：II.3.9}
\NormalTok{多变量的：R.3.13}
\NormalTok{两个变元的：II.3.9}
\NormalTok{函数图像：II.3.4}
\NormalTok{函数关系：I.5.3，R.2.1}
\NormalTok{函数符号：I.5.3}

\NormalTok{通解：I.5.2}
\NormalTok{（序列的）通项：R.7.8}
\NormalTok{广义连续统假设：III.6.4}
\NormalTok{（集合的）一般元素：R.1.2}
\NormalTok{一般结构：IV.1.4}
\NormalTok{图像：II.3.1}
\NormalTok{函数的：II.3.4}
\NormalTok{逆的：II.3.2}
\NormalTok{对应的：II.3.1}
\NormalTok{映射的：R.3.5}
\NormalTok{关系的：II.3.1，R.3.2}
\end{Highlighting}
\end{Shaded}

对称图像：II.3.2 图像的复合：II.3.3 大于：III.1.3，R.6.3 有序集的最大元素：III.1.7，R.6.5 最大下界：III.1.9，R.6.7

半开区间：III.1.13，R.6.4 假设，辅助的：I.3.3 连续统：III.6.4 广义连续统：III.6.4 归纳的：III.2.2，III.4.3

等同：R.8.5 恒等：R.1.3 恒等对应：II.3.3 恒等映射：II.3.4，R.2.3 像，直接（结构的）：IV.2.6 逆（集合在映射下的）：R.2.6 逆（结构的）：IV.2.4 覆盖的：II.4.6 等价关系的：II.6.6 函数（滥用语言）的：R.2.4 集合的：II.3.2 集合在对应下的：II.3.1 集合在函数下的：II.3.4，R.2.4 集合在图像下的：II.3.1 隐公理：I.2.1 蕴含：R.1.3 不可达基数：III.6.Ex.22 不可达初始序数：III.6.Ex.16 包含关系：II.1.2，III.1.1，R.1.2 不相容公理：R.8.6 子集的递增族：III.1.5，R.6.12 递增映射：III.1.5，R.6.12 不可分解序数：III.2.Ex.16 独立（于其他公理）：I.2.Ex.1 族的指标集：II.3.4，R.2.2 指标记号：R.2.2 诱导等价关系：R.5.5 诱导映射：II.6.5 诱导序、序关系、预序、预序关系：III.1.4，R.6.1 诱导关系：II.6.3 诱导结构：IV.2.4

归纳，降：III.4.3 诺特的：III.6.5 原理：III.2.2，III.4.3 受限的：III.4.3 从k开始：III.4.3 超限：III.2.2 归纳假设：III.2.2，III.4.3 归纳集：III.2.4，R.6.9 下确界：III.1.9，R.6.7 无穷序列：III.6.1，R.7.8 无穷集：III.6.1 无穷公理：III.6.1 初始序数：III.6.Ex.10 词的初始段：I.App.1 初始结构：IV.2.3 单射：II.3.7，R.2.8 典范的：II.3.7 单射映射：II.3.7，R.2.8 整数，自然数：III.4.1 两个整数之商的整数部分：III.5.6 相交的子集：R.1.13 交集，可数的：R.7.9 集合族或子集族的：II.4.1，R.4.6 集合族的：II.4.1 若干集合的：R.1.13 区间（闭、半开、开、无界）：III.1.13，R.6.4 非传递的：II.6.Ex.11 固有项：IV.1.6 不变元素：R.2.3 映射的逆扩张：R.2.6 逆，（双射）映射的：II.3.7，R.2.9 对应的：II.3.2 图像的：II.3.2 逆像：R.2.6 覆盖的：II.4.6 等价关系的：II.6.6 集合的：II.3.2 结构的：IV.2.4 逆同构：IV.1.5 逆向极限，映射族的：III.7.2，R.6.14 集合族的：III.7.1，R.6.14 逆向系统，映射的：III.7.2，R.6.14 集合的：III.7.1，R.6.14 子集的：III.7.2

对合置换：II.3.7，R.2.9 不可约元素（格的）：III.4.Ex.7 同构的集合：IV.1.5 结构：IV.1.5，R.8.6 同构：IV.1.5，R.8.6 有序集的：III.1.3 映射的迭代：III.6.2，R.2.11

大于：III.1.3 有限序列的末项：III.5.4 格：III.1.11，R.6.8 布尔：III.1.Ex.17 完备：III.1.Ex.11 分配：III.1.Ex.16 相对有补：III.1.Ex.17 有序集的最小元素：III.1.7，R.6.5 最小上界：III.1.9，R.6.7 基数族的：III.3.2 左有向集：III.1.10，R.6.8 区间的左端点：III.1.13 左逆（单射的）：II.3.8 合法化定理：I.3.3 引理：I.2.2 佐恩：III.2.4，R.6.10 长度，有限序列的：III.5.4 词的：I.App.1 小于：III.1.3，R.6.3 字典序关系、序：III.2.6 积：III.2.6 和 III.2.Ex.10 极限，正向：III.7.5，R.6.13 逆向：III.7.1，R.6.14 线性有序：见全序 链：I.1.1 关系的逻辑成分：I.App.Ex.6 逻辑构造：I.App.Ex.6 逻辑符号：I.1.1 逻辑理论：I.3.1 逻辑构造的：I.App.Ex.6 逻辑不可约的：I.App.Ex.6 下界：III.1.8，R.6.7 最大：III.1.9，R.6.7

映射：R.2.1 双射的：II.3.7，R.2.9

\begin{Shaded}
\begin{Highlighting}[]
\NormalTok{上有界（下有界，有界）：III.1.8，R.6.7}
\NormalTok{与等价关系相容：II.6.5}
\NormalTok{与两个等价关系相容：R.5.8}
\NormalTok{复合：II.3.7，R.2.11}
\NormalTok{常数：II.3.4，R.2.3}
\NormalTok{递减：III.1.5，R.6.12}
\NormalTok{对角线：II.3.7，II.5.3，R.3.4}
\NormalTok{空集：II.3.4}
\NormalTok{恒等式：II.3.4，R.2.3}
\NormalTok{递增：III.1.5，R.6.12}
\NormalTok{单射：II.3.7，R.2.8}
\NormalTok{逆映射：II.3.7，R.2.9}
\NormalTok{单调性：III.1.5}
\NormalTok{将一个集合映入另一个集合的映射：II.3.4}
\NormalTok{一个集合到另一个集合上的映射：II.3.7}
\NormalTok{一{-}对{-}一：II.3.7，R.2.9}
\NormalTok{阶{-}保持：III.1.5}
\NormalTok{阶{-}反序：III.1.5}
\NormalTok{偏序：II.3.9，R.3.13}
\NormalTok{严格递减（严格递增，严格单调）：III.1.5，R.6.12}
\NormalTok{满射：II.3.7，R.2.4}
\NormalTok{泛的：IV.3.1}
\NormalTok{映射，在集合上一致：II.3.5}
\NormalTok{典范的：见“典范”}
\NormalTok{数学理论：I.1.1, I.2.1, I.2.2}
\NormalTok{极大元：II.1.6, R.6.7}
\NormalTok{最大值：R.6.5}
\NormalTok{属于，关系：II.1.1，R.1.10}
\NormalTok{方法，分情形讨论：I.3.3}
\NormalTok{归谬法的：I.3.3}
\NormalTok{辅助常数的：I.3.3}
\NormalTok{辅助假设的：I.3.3}
\NormalTok{极小元：III.1.6，R.6.6}
\NormalTok{最小值：R.6.5}
\NormalTok{移动的（有限子集的集合）：III.4.Ex.11}
\NormalTok{理论模型：I.2.4}
\NormalTok{单调映射：III.1.5}
\NormalTok{态射：IV.2.1}
\NormalTok{整数的倍数：III.5.6}
\NormalTok{多重序列：III.6.1}
\NormalTok{多值理论：R.8.7}
\NormalTok{互不相交的（集合族）：II.4.7}
\NormalTok{自然整数：III.4.1}
\end{Highlighting}
\end{Shaded}

\begin{Shaded}
\begin{Highlighting}[]
\NormalTok{关系的否定：I.1.3}
\NormalTok{诺特归纳原理：III.6.5}
\NormalTok{诺特集：III.6.5}
\NormalTok{映射的第n次迭代：III.6.2，R.2.11}
\NormalTok{序列的第n项：III.6.1，R.7.8}
\NormalTok{有限集的元素个数：III.4.1}
\NormalTok{数字符号：III.5.7}
\NormalTok{记数系统：III.5.7}

\NormalTok{奇数整数：III.5.6}
\NormalTok{一{-}对{-}一对应，映射：II.3.7，R.2.9}
\NormalTok{开区间：III.1.13，R.6.4}
\NormalTok{相反序关系、预序关系：III.1.1，III.1.2，R.6.1}
\NormalTok{有序对：II.2.1}
\NormalTok{序{-}保持映射：III.1.5}
\NormalTok{与预序关系相关的序关系：III.1.2}
\NormalTok{在x与y之间，关于x与y：III.1.1}
\NormalTok{在集合上诱导的：III.1.4}
\NormalTok{字典序的：III.2.6}
\NormalTok{在集合上：III.1.1，R.6.1}
\NormalTok{与序关系相反的：III.1.1}
\NormalTok{序{-}反转映射：III.1.5}
\NormalTok{序{-}结构：R.6.1}
\NormalTok{总数：III.1.12}
\NormalTok{序型：III.2.习题13}
\NormalTok{有序集：III.1.3，R.6.1}
\NormalTok{反定向的：III.1.习题23}
\NormalTok{分支的：III.1.习题24}
\NormalTok{完全分歧的：III.2.习题8}
\NormalTok{部分良{-}序的：III.2.习题4}
\NormalTok{分歧的：III.2.习题8}
\NormalTok{分散的：III.1.习题20}
\NormalTok{无间隙的：III.1.习题19}
\NormalTok{序：III.1.1，R.6.1}
\NormalTok{诱导的：III.1.4}
\NormalTok{字典序的：III.2.6}
\NormalTok{积：III.1.4}
\NormalTok{总数：III.1.12}
\NormalTok{序数：III.2.习题14}
\NormalTok{临界的：III.6.习题13}
\NormalTok{函数符号：III.2.习题17}
\NormalTok{不可达初始：III.6.习题16}
\NormalTok{不可分解的：III.2.习题16}
\NormalTok{序数，初始：III.6.习题10}
\end{Highlighting}
\end{Shaded}

序数积（序型的）：III.2.习题13 正则的：III.6.习题16 奇异的：III.6.习题16 序数和，由有序集索引的非空有序集族的：III.1.习题3（序型的）：III.2.习题13

\begin{Shaded}
\begin{Highlighting}[]
\NormalTok{有序对：II.2.1}
\NormalTok{两两不相交（子集族）：R.4.4}
\NormalTok{参数，参数集，参数表示：II.3.7，R.2.14}
\NormalTok{部分映射：II.3.9，R.3.13}
\NormalTok{部分积：II.5.4}
\NormalTok{部分良{-}序的（有序集）：III.2.习题1}
\NormalTok{集合的划分：II.4.7，R.4.4}
\NormalTok{到商集的过渡：II.6.3，II.6.5，R.5.7，R5.8}
\NormalTok{完全平衡的组装：I.附录4}
\NormalTok{置换：II.3.7，R.2.9}
\NormalTok{较贫（结构物种）：IV.1.6}
\NormalTok{幂，集合的：III.3.1，R.7.2}
\NormalTok{连续统的幂：III.6.4}
\NormalTok{幂，等价的：R.7.2}
\NormalTok{幂的和：R.7.5}
\NormalTok{前驱（序数的）：III.2.Ex.14}
\NormalTok{预序关系：III.1.2，R.6.1}
\NormalTok{集合上的预序关系：III.1.2}
\NormalTok{与预序关系相反的关系：III.1.2}
\NormalTok{预序集：R.6.1}
\NormalTok{预序：III.1.2}
\NormalTok{（较粗，较细）：III.1.4}
\NormalTok{主基集：IV.1.3，1.4}
\NormalTok{归纳原理：III.4.3}
\NormalTok{诺特归纳原理：III.6.5}
\NormalTok{超限归纳的：III.2.2}
\NormalTok{结构的演绎程序：IV.1.6}
\NormalTok{乘积基数：III.3.3}
\NormalTok{乘积，字典序：III.2.6}
\NormalTok{一族映射的积：II.5.7}
\NormalTok{一族集合的积：II.5.3，R.4.9}
\NormalTok{基数的积：III.3.3}
\NormalTok{序关系、序、预序关系、预序的积：III.1.4}
\NormalTok{有序集、预序集的：III.1.4}
\NormalTok{多个集合的：R.3.12}
\NormalTok{关于两个覆盖：II.4.6}
\NormalTok{两个等价关系的：II.6.8，R.5.10}
\end{Highlighting}
\end{Shaded}

\begin{Shaded}
\begin{Highlighting}[]
\NormalTok{两个集合的积：II.2.2，R.3.1}
\NormalTok{序数积：III.2.Ex.13}
\NormalTok{积结构，结构的积：IV.2.4}
\NormalTok{偏积：II.5.4}
\NormalTok{投影，图的（第一，第二）坐标：II.3.1}
\NormalTok{有序对的（第一，第二）坐标：II.2.1}
\NormalTok{到因子上的投影：II.5.3}
\NormalTok{到部分积上的投影：II.5.4}
\NormalTok{投影：R.3.1, R.3.12, R.4.11}
\NormalTok{证明：I.2.2}
\NormalTok{真段：I.App.1}
\NormalTok{有限特征性质：III.4.5}
\NormalTok{命题：I.2.2}
\NormalTok{伪{-}序数：III.2.Ex.20}

\NormalTok{量化理论：I.4.2}
\NormalTok{量词，存在：I.4.1}
\NormalTok{典型的：I.4.4}
\NormalTok{全称：I.4.1}
\NormalTok{商等价关系：II.6.7, R.5.9}
\NormalTok{预序关系或预序集关于等价关系的商：III.1.Ex.2}
\NormalTok{两个整数的商：III.5.6}
\NormalTok{商集：II.6.2, R.5.2}
\NormalTok{商结构：IV.2.6}

\NormalTok{分叉的（有序集）：III.2.Ex.8}
\NormalTok{值域，对应的：II.3.1}
\NormalTok{图的：II.3.1}
\NormalTok{实现，阶梯类型的：IV.1.Ex.1}
\NormalTok{带符号阶梯类型的：IV.2.Ex.1}
\NormalTok{归谬法，方法：I.3.3}
\NormalTok{覆盖的加细：II.4.6}
\NormalTok{集合上的自反关系：II.6.1, R.5.1}
\NormalTok{正则基数：III.6.Ex.17}
\NormalTok{正则序数：IV.6.Ex.16}
\NormalTok{关系：I.1.3}
\NormalTok{A的元素与B的元素之间的关系：II.3.1}
\NormalTok{集合化的：II.1.4}
\NormalTok{与等价关系相容的：II.6.3, R.5.7}
\NormalTok{与两个等价关系相容的：II.6.8}
\NormalTok{等价：II.6.1, R.5.2}
\NormalTok{假：I.2.2}
\NormalTok{函数的：I.5.3, R.2.1}
\end{Highlighting}
\end{Shaded}

\begin{Shaded}
\begin{Highlighting}[]
\NormalTok{关系，具有图的：II.3.1}
\NormalTok{    由等价关系在过渡到商时诱导的：II.6.3}
\NormalTok{    相等：I.5.1, R.1.6}
\NormalTok{    包含：II.1.2, III.1.1, R.1.12}
\NormalTok{    属于：II.1.1, R.1.10}
\NormalTok{    序：III.1.1, R.6.1}
\NormalTok{    预序：III.1.2, R.6.1}
\NormalTok{    自反的：II.6.1, R.5.1}
\NormalTok{    单{-}值的：I.5.3}
\NormalTok{    对称的：II.6.1, R.3.4}
\NormalTok{    全序：II.1.12}
\NormalTok{    传递的：II.6.1, R.5.1}
\NormalTok{    可迁移的：IV.1.3}
\NormalTok{    真：I.2.2}
\NormalTok{    良{-}排序：III.2.1}
\NormalTok{    ≤（相应地 ≥，⊂，⊃）：III.1.3}
\NormalTok{关系，等价的：I.3.5, R.1.3}
\NormalTok{关系符号：I.1.3}
\NormalTok{相对补（格中元素的）：III.1.Ex.17}
\NormalTok{相对有补格：III.1.Ex.17}
\NormalTok{余数（a除以b的）：III.5.6}
\NormalTok{表示，参数：II.3.7, R.2.14}
\NormalTok{代表，等价类的：II.6.2}
\NormalTok{代表，系统：II.6.2}
\NormalTok{受限归纳：III.4.3}
\NormalTok{直接系的限制：III.7.6}
\NormalTok{    函数的限制：II.3.5，R.2.13}
\NormalTok{    逆向系的限制：III.7.1}
\NormalTok{单射的收缩：II.3.8}
\NormalTok{更丰富的结构：R.8.3}
\NormalTok{右有向集：III.1.10，R.6.8}
\NormalTok{右{-}区间的右端点：III.1.13}
\NormalTok{（满射的）右逆：II.3.8}

\NormalTok{饱和集：R.5.6}
\NormalTok{集合关于等价关系的饱和：R.5.6}
\NormalTok{集合的标度：R.8.1}
\NormalTok{散布的（有序集）：III.1.Ex.20}
\NormalTok{方案：I.2.1}
\NormalTok{    分层构造：IV.1.1}
\NormalTok{    选择与并的构造：II.1.6}
\NormalTok{有序对的第二坐标：II.2.1}
\NormalTok{    第二投影：II.3.1}
\NormalTok{    第二类：I.1.3}
\end{Highlighting}
\end{Shaded}

\begin{Shaded}
\begin{Highlighting}[]
\NormalTok{对应关系的截面：II.3.1}
\NormalTok{图的：II.3.1}
\NormalTok{（关于等价关系的）集合的截面：II.6.2}
\NormalTok{乘积的子集的截面：R.3.7}
\NormalTok{满射的截面：II.3.8}
\NormalTok{段，终段：I.App.1}
\NormalTok{初段：I.App.1}
\NormalTok{词的段：I.App.1}
\NormalTok{汇编的段：I.App.4}
\NormalTok{真段：I.App.1}
\NormalTok{以x为端点的段：III.2.1}
\NormalTok{段，不相交的：I.App.1}
\NormalTok{选择与并集，方案：II.1.6}
\NormalTok{分离（E 的元素被 α 分离）{-}映射）：IV.3.3}
\NormalTok{序列：R.7.8}
\NormalTok{双重：R.7.8}
\NormalTok{有限：III.5.4，R.7.8}
\NormalTok{无限：III.6.1，R.7.8}
\NormalTok{通过由映射定义的顺序排列可数族而获得：III.6.1}
\NormalTok{集合的元素：III.6.1}
\NormalTok{显著：I.App.2}
\NormalTok{平稳的：III.6.5}
\NormalTok{三重，多重：III.6.1}
\NormalTok{仅项的顺序不同的序列：III.6.1，R.7.8}
\NormalTok{集合：II.1.1}
\NormalTok{基：IV.1.3，1.4}
\NormalTok{有界的，上有界的，下有界的：III.1.8，R.6.7}
\NormalTok{由单个元素组成：II.1.5}
\NormalTok{可数的：III.6.4，R.7.7}
\NormalTok{适当的：III.2.Ex.20}
\NormalTok{有向的：III.1.10，R.6.8}
\NormalTok{空的：II.1.7，R.1.8}
\NormalTok{赋予结构的：IV.1.4}
\NormalTok{与另一集合等势的：R.7.1}
\NormalTok{有限的：III.4.1}
\NormalTok{（族的）指标：II.3.4}
\NormalTok{归纳的：III.2.4，R.6.9}
\NormalTok{无限的：III.6.1}
\NormalTok{左有向的：III.1.10，R.6.8}
\NormalTok{诺特的：III.6.5}
\NormalTok{所有满足P的x∈A：II.1.6}
\NormalTok{所有满足R的x（R为在x中收集化的关系）：II.1.4}
\NormalTok{基数≤a的：III.3.2}
\end{Highlighting}
\end{Shaded}

\begin{Shaded}
\begin{Highlighting}[]
\NormalTok{等价对象类之集合：II.5.9}
\NormalTok{族中元素之集合：II.3.4，R.2.14}
\NormalTok{从E到F的映射之集合：II.5.2，R.2.2}
\NormalTok{n个元素之集合（n为整数）：III.4.1}
\NormalTok{形如T（x∈A）的对象之集合：II.1.6}
\NormalTok{参数表示之参数集：II.3.7}
\NormalTok{σ{-}态射之集合：IV.2.1}
\NormalTok{集合的子集之集合：II.5.1，R.1.10}
\NormalTok{有限特征子集之集合：III.4.5}
\NormalTok{有序的：III.1.3，R.6.1}
\NormalTok{预序的：III.1.3，R.6.1}
\NormalTok{积：R.3.1，R.3.12，R.4.9}
\NormalTok{商：II.6.2，R.5.2}
\NormalTok{代表（关系的）：II.3.1}
\NormalTok{右有向的：III.1.10，R.6.8}
\NormalTok{全序的：III.1.12，R.6.4}
\NormalTok{传递的：III.2.习题20}
\NormalTok{泛的：IV.3.1}
\NormalTok{良{-}序的：III.2.1，R.6.5}
\NormalTok{唯一元素为x的：II.1.5}
\NormalTok{集合，复合：R.3.10}
\NormalTok{不相交的：II.4.7}
\NormalTok{等势的：III.3.1}
\NormalTok{同构的：IV.1.5}
\NormalTok{理论：II.1.1}
\NormalTok{Σ{-}容许的：IV.3.2}
\NormalTok{σ{-}态射：IV.2.1}
\NormalTok{符号：I.1.1，I.附录1}
\NormalTok{逻辑的：I.1.1}
\NormalTok{有理的：I.1.3}
\NormalTok{特定的：I.1.1}
\NormalTok{实义的：I.1.3}
\NormalTok{带符号的典范扩张：IV.2.习题1}
\NormalTok{带符号的阶梯型：IV.2.习题1}
\NormalTok{有意义的序列、词：I.附录2}
\NormalTok{单{-}值关系：I.5.3}
\NormalTok{奇异基数：III.6.习题17}
\NormalTok{奇异序数：III.6.习题16}
\NormalTok{小于：IV.1.3}
\NormalTok{解（方程的解）：I.2.2}
\NormalTok{完全的（或一般的）：I.5.2}
\NormalTok{泛映射问题的解：IV.3.1}
\NormalTok{对应的源：II.3.1}
\NormalTok{种，等价的：IV.1.7}
\end{Highlighting}
\end{Shaded}

\begin{Shaded}
\begin{Highlighting}[]
\NormalTok{代数结构之种：IV.1.4}
\NormalTok{序结构之种：IV.1.4}
\NormalTok{结构之种：IV.1.4，R.8.2}
\NormalTok{拓扑结构之种：IV.1.4}
\NormalTok{较贫的（较富的）：IV.1.6}
\NormalTok{特定符号：I.1.1}
\NormalTok{稳定子集：R.2.5}
\NormalTok{平稳序列：III.6.5}
\NormalTok{严格上界：III.2.4}
\NormalTok{严格较粗（较细）的结构：IV.2.2}
\NormalTok{严格递减（递增）的：III.1.5，R.6.12}
\NormalTok{严格大于（小于）：III.1.3，R.6.3，R.7.3}
\NormalTok{较强的理论：I.2.4}
\NormalTok{强不可达基数：III.6.习题22}
\NormalTok{结构，较粗的：IV.2.2}
\NormalTok{导出的：IV.1.6}
\NormalTok{最终的：IV.2.5}
\NormalTok{更细的：IV.2.2}
\NormalTok{泛型：IV.1.4}
\NormalTok{诱导：IV.2.4}
\NormalTok{初始：IV.2.3}
\NormalTok{集合上的：IV.1.4, R.8.2}
\NormalTok{物种Σ：IV.1.4}
\NormalTok{阶：R.6.1}
\NormalTok{积：IV.2.4}
\NormalTok{商：IV.2.6}
\NormalTok{更丰富：R.8.2}
\NormalTok{严格更粗（更细）：IV.2.2}
\NormalTok{从属：IV.1.6}
\NormalTok{的传输：R.8.5}
\NormalTok{底层：IV.1.6}
\NormalTok{结构，可比较的：IV.2.2}
\NormalTok{同构的：IV.1.5, R.8.6}
\NormalTok{子族：II.3.5，R.2.14}
\NormalTok{子格：III.4.Ex.9}
\NormalTok{从属结构：IV.1.6}
\NormalTok{子序列（序列的）：III.6.1，R.8.7}
\NormalTok{子集，共尾的：III.1.7，R.6.5}
\NormalTok{共初始的：III.1.7，R.6.5}
\NormalTok{仅由a组成的：R.1.9}
\NormalTok{空的：R.1.8}
\NormalTok{集合的：II.2.1，R.1.7}
\NormalTok{饱和的（关于等价关系）：II.6.4}
\NormalTok{对称的：R.3.4}
\end{Highlighting}
\end{Shaded}

子集，不相交的：R.1.13 集合：II.5.1，R.1.10 稳定的：R.2.5 实体化符号：I.1.3 代入，判定准则：I.1.2 和，基数的：III.3.3 集合族的：II.4.8 基数的：III.3.3 幂的：R.7.5 集合的：R.4.5 序数的：III.2.Ex.13 上确界：III.1.9，R.6.7 满射：II.3.7，R.2.4 满射映射：II.3.7，R.2.4 符号，函数的：I.5.3 数值的：III.5.7 对称图：II.2.3 关系或子集：II.6.1，R.3.4 对称性，典范的：R.3.4 系统，十进制的：III.5.7 直（直接）的：III.7.5，R.6.13 二进制的：III.5.7 逆（反向）的：III.7.1，R.6.14 记数系统：III.5.7 代表系：II.6.2

对应的目标：II.3.1 项：I.1.3 一般的：R.7.8 （有限序列的）（第一、第k、最后）项：III.5.4 内在的：IV.1.6 第n项：R.7.8 满足关系的：I.2.2 可写成T形式的：I.5.2 文本，论证性的：I.2.2 定理：I.2.2 康托尔的：III.3.6 合法化的：I.3.3 策梅洛的：III.2.3，R.6.5 理论，等价的：I.2.4 理论：I.1.1，I.2.1，I.2.2 矛盾的：I.2.2 等式的：I.5.1 逻辑的：I.3.1

\begin{Shaded}
\begin{Highlighting}[]
\NormalTok{理论 多值的：R.8.7}
\NormalTok{结构类的：IV.1.4}
\NormalTok{集合的：II.1.1}
\NormalTok{给定物种的结构：R.8.2}
\NormalTok{量化：I.4.2}
\NormalTok{更强：I.2.4}
\NormalTok{单叶：R.8.7}
\NormalTok{全序关系，全序：III.1.12}
\NormalTok{全序集：III.1.12，R.6.4}
\NormalTok{集合族的迹：II.4.5}
\NormalTok{子集或子集族的迹：R.1.16}
\NormalTok{元素在函数下的变换：II.3.4，R.2.4}
\NormalTok{传递关系：II.6.1，R.5.1}
\NormalTok{传递集：III.2.Ex.20}
\NormalTok{传递性判据：IV.2.3，2.5}
\NormalTok{结构的传递：R.8.5}
\NormalTok{可传递关系：IV.1.3}
\NormalTok{传输结构：IV.1.5}
\NormalTok{横截的：II.6.2}
\NormalTok{三重：II.2.2}
\NormalTok{三重序列：III.6.1}
\NormalTok{真关系：I.2.2}
\NormalTok{结构种类的一个典型刻画：IV.1.4}
\NormalTok{典型量词：I.4.4}
\NormalTok{典型化：IV.1.3}

\NormalTok{无界区间：III.1.13，R.6.4}
\NormalTok{底层结构：IV.1.6}
\NormalTok{并，可数：R.7.9}
\NormalTok{集合族的并：R.4.2}
\NormalTok{集合的集合的并：II.4.1}
\NormalTok{若干集合的并：R.1.13}
\NormalTok{单值结构物种：IV.1.5}
\NormalTok{单值理论：R.8.7}
\NormalTok{泛映射：IV.3.1}
\NormalTok{泛问题：IV.3.1}
\NormalTok{全称量词：I.4.1}
\NormalTok{全集：IV.3.1}
\NormalTok{上界：III.1.8，R.6.7}
\NormalTok{最少：III.1.9，R.6.7}
\NormalTok{严格：III.2.4}

\NormalTok{函数的值：II.3.4}
\NormalTok{函数在某元素处的值：R.2.1}
\end{Highlighting}
\end{Shaded}

变量的值：R.1.2 由对应关系取得：II.3.1 变量：R.1.2 总体的方差：IV.2.Ex.1

弱相容性（等价关系与预序关系的）：III.1.Ex.2

策梅洛公理（=选择公理）：R.4.9

策梅洛定理：III.2.3，R.6.5

佐恩引理：III.2.4，R.6.10

